Ask a biomedical search engine ``What is the function of DITPA?'' and it retrieves dozens of passages about thyroid hormone analogs. A clinician scanning those results spots the connection immediately. DITPA is a thyroid hormone analog. Its mechanism sits under the Medical Subject Headings (MeSH) tree node for thyroid hormones. A standard retrieval pipeline scores each passage against the query by lexical overlap or embedding similarity, and it places the most relevant evidence at rank 54. The passage is sitting in the candidate pool. Pairwise scoring cannot see it.

Across 943 held-out BioASQ test questions, a strong Hybrid RRF retriever leaves every gold evidence passage outside the top 10 for over 6\% of queries where that evidence does exist somewhere in the top 100. This happens because biomedical evidence lives in higher-order relations. Questions, passages, MeSH concepts, biomedical entities, documents, and knowledge-base entries connect through pathways that pairwise query-passage comparisons cannot trace. A passage about DITPA and a passage about levothyroxine share no surface lexical overlap. They sit far apart in embedding space. They belong to the same MeSH subtree, and a pairwise scoring function never asks about that.

Existing biomedical retrievers~\citep{jin2023medcpt} and medical RAG systems~\citep{xiong2024benchmarking,xiong2024imedrag} score each passage independently against the query. They do not ask whether two candidate passages share a MeSH ancestor. They do not check whether the biomedical entities in a question form a dense cluster with those of a low-ranked passage. GraphRAG-style approaches~\citep{edge2024graphrag,wu2024medgraphrag,luo2025hypergraphrag} do apply graph structure to retrieval, but their edges come from lexical overlap or embedding similarity rather than domain-grounded concept annotations such as MeSH, biomedical entity linking, or curated knowledge graphs like PrimeKG. No existing work isolates whether raw graph topology, stripped of medical knowledge, helps or actively hurts biomedical evidence ranking.

We introduce KCH-MedRank, a knowledge-constrained learning-to-rank framework for biomedical evidence reranking. The method starts from a reproducible candidate pool built with BM25, dense retrieval, and Reciprocal Rank Fusion (RRF). It enriches each query's local candidate neighborhood with four categories of features: (1) retrieval scores from lexical, dense, and fused sources; (2) biomedical semantic features from MeSH hierarchy annotations and entity overlap metrics; (3) relation features from the PrimeKG biomedical knowledge graph; and (4) optional structural features from a local hypergraph that captures n-ary concept interactions that pairwise graphs cannot preserve. All features are fed into query-grouped LambdaMART---a gradient-boosted tree-based learning-to-rank algorithm~\citep{burges2010lambdamart}---yielding a transparent, feature-attributable reranker. The framework is designed as a focused evidence selection module that can serve as a retrieval stage in medical RAG pipelines, though the experiments in this paper center on retrieval and reranking metrics.

As controlled baselines we use a flat knowledge-feature model that removes the hypergraph entirely, and a pairwise graph decomposition that breaks multi-node hyperedges into anchor-based edges. The core experimental design isolates three questions. \textbf{RQ1:} Do flat biomedical knowledge features alone improve on retrieval signals? \textbf{RQ2:} Does hypergraph structure provide additional benefit, and under what conditions? \textbf{RQ3:} Are domain-grounded concept annotations---MeSH, entity, PrimeKG---necessary for graph-based biomedical reranking, or would generic graph topology built from lexical or embedding similarity suffice?

Our contributions are the following. First, a knowledge-constrained learning-to-rank framework that integrates retrieval, biomedical semantic, MeSH hierarchy, entity, relation, and optional hypergraph features for interpretable evidence reranking. Second, a controlled ablation showing that flat biomedical knowledge features explain most aggregate gains; hypergraph structure provides targeted but modest early-rank benefit concentrated in questions lacking direct MeSH overlap and in multi-evidence scenarios; and removing domain-grounded concept annotations from the hypergraph degrades ranking, suggesting that generic graph topology built without medical knowledge is insufficient for productive graph-based reranking. Third, on BioASQ, statistically significant Recall@10 improvement over a true MedCPT Cross-Encoder baseline ($+0.0157$, $p=0.0076$) with $18.6\times$ faster reranking (reranking-stage latency under cached features, excluding candidate generation), plus supplementary diagnostics on PubMedQA and BEIR NFCorpus. Fourth, a diagnostic analysis categorizing 648 of 661 rescued gold passages by dominant activated knowledge mechanism, with MeSH hierarchy paths accounting for 75.9\% of rescues, shared entity clusters for 23.9\%, and PrimeKG relations for the remainder.

We position this work as a focused evidence selection module for medical RAG rather than a complete clinical QA system. The method does not replace clinical judgment or guarantee answer correctness. The failure analysis reveals cases where gold passages remain beyond the reach of our knowledge features, and the hypergraph benefit over flat knowledge features is modest and concentrated at early ranks ($k \leq 5$). These limitations are discussed transparently, along with directions for broader knowledge-graph coverage and generation-level evaluation.
